% ЕМТ 2024, XI клас — точен LaTeX препис от официалните файлове в Klasirane.com.
% Условия: https://klasirane.com/api/competitions/EMT/All/2024/11%20%D0%BA%D0%BB/probs
% SHA-256 (условия): 231e1da6860db8c04ee04f0d21609fca6d4bad4863d83ffaae0e2670b3767f8f
% Решения: https://klasirane.com/api/competitions/EMT/All/2024/11%20%D0%BA%D0%BB/sol
% SHA-256 (решения): 4e3dcb172fa7519ad4bef320cf3562e535a957ac053b0a8ceab2db2001894fc1
% Точковите оценки, правописът и формулировките са запазени от източника.
% Файловете са сканирани изображения; всички видими страници са проверени пряко.
% Първата страница на файла с решенията е заглавна и не съдържа условия или решения.

Задача 11.1. Да се намерят всички реални числа $a$, за които уравнението

$$x^3-(a+2)x^2-(a-2)x+2a-1=0$$

има три различни корена $x_1$, $x_2$ и $x_3$, и тези корени заедно с числото $a$ в някакъв ред образуват аритметична прогресия.

Решение. Тъй като

$$x^3-(a+2)x^2-(a-2)x+2a-1=(x-1)(x^2-(a+1)x-2a+1),$$

то $x_1=1$, а $x_2$ и $x_3$ са корени на $f(x)=x^2-(a+1)x-2a+1=0$. Тъй като $x_2+x_3=a+1$, то за дадената аритметична прогресия, с точност до симетрия, има две възможности – $x_2,1,a,x_3$ или $1,x_2,x_3,a$.

В първия случай $x_2=2-a$ и след заместване в уравнението получаваме

$$f(2-a)=0\Longleftrightarrow2a^2-7a+3=0$$

с корени $a=3$ и $a=\frac12$.

Във втория случай $d=x_2-1$, като $x_2=\frac{a+2}{3}$ и получаваме

$$f\left(\frac{a+2}{3}\right)=0\Longleftrightarrow2a^2+23a-7=0,$$

с корени $a=\frac{-23+\sqrt{585}}4$ и $a=\frac{-23-\sqrt{585}}4$.

Оценяване. (6 точки) 1 т. за разлагането $(x-1)(x^2-(a+1)x-2a+1)$; 1 т. за наблюдението, че $x_2+x_3=a+1$ води до разглеждането на само два случая; по 2 т. за пълно решаване на всеки от двата случая.

Задача 11.2. Ъглите при върховете $A$, $B$ и $C$ на триъгълник $ABC$ са съответно първи, втори и трети член на намаляваща аритметична прогресия. Намерете ъглите на триъгълника, ако $\angle BHI=60^\circ$, където $H$ и $I$ са съответно ортоцентърът и центърът на вписаната окръжност на триъгълника.

Отговор. $80^\circ,60^\circ,40^\circ$.

Решение. Първи метод. От условието следва, че $\angle B=60^\circ$. Тогава $\angle AIC=120^\circ$ и ако $P$ е симетричната на $I$ спрямо $AC$, то $\angle APC=120^\circ$. Това означава, че $P$ е върху описаната около $ABC$ окръжност. Тъй като симетричната на $H$ спрямо $AC$ също лежи на описаната окръжност (означаваме тази точка с $Q$), то $QHIP$ е равнобедрен трапец. Следователно $\angle PQH=\angle IHQ=120^\circ$.

Четириъгълникът $QBCP$ е вписан в окръжност, откъдето получаваме $\angle PCB=180^\circ-\angle BQP=60^\circ$. Понеже $\angle BCI=\angle ICA=\angle ACP$, получаваме $\angle ACB=40^\circ$ и $\angle BAC=80^\circ$.

Втори метод. (Б. Димитров) Ще използваме стандартно означение за ъглите на триъгълника $ABC$. Тогава от даденото условие получаваме $2\beta=\alpha+\gamma\Longrightarrow3\beta=180^\circ\Longrightarrow\beta=60^\circ$.

Ще решим задачата за остроъгълен триъгълник (когато $\triangle ABC$ е тъпоъгълен, разсъжденията са аналогични). Имаме $\angle AHC=\angle AIC=120^\circ$, откъдето четириъгълникът $AHIC$ е вписан. Сега

$$180-\alpha=\angle BHC=\angle BHI+\angle IHC=60^\circ+\angle IAC=60^\circ+\frac\alpha2\Longrightarrow\alpha=80^\circ$$

От последното получаваме $\gamma=40^\circ$.

Оценяване. (6 точки) Първи метод: 1 т. за $\beta=60^\circ$; 1 т. за доказване, че симетричната на $I$ лежи на описаната окръжност; 2 т. за вписания четириъгълник $QBCP$; 2 т. за намиране на ъглите. Втори метод: 1 т. за $\beta=60^\circ$; 2 т. за вписания четириъгълник $AHIC$; 3 т. за намиране на ъглите.

Задача 11.3. Кристи иска да раздаде бонбони на $n\geq3$ свои съученици. Той разполага всеки от тях върху точка в двора на училището. Точките са в една равнина, като никои три от тях не лежат на една права. За всеки изпъкнал многоъгълник $P$ с върхове сред тези точки, Кристи прави следното. Преброява учениците, които се намират вътре в $P$ или на страните му, нека техният брой е $S_P$. Той раздава по $S_P$ бонбона на всеки от тези $S_P$ ученици. Ученик, получил най-малко бонбони след всички раздавания, наричаме нещастен (нещастните ученици могат да са един или повече). Определете максималното количество бонбони, които може да получи нещастен ученик.

Решение. Нека $X$ е множеството от $n$ точки. За изпъкнал многоъгълник $P$ с върхове в $X$ да означим с $C(P)$ множеството от точки в $X$, които са вътре или на контура на $P$. Нека $Q$ е общият брой бонбони, получени от всички ученици, а $m$ е броят бонбони, получени от нещастен ученик. Имаме

$$Q=\sum_P|C(P)|^2\leq\sum_{A\subset X}|A|^2=\sum_{k=3}^nk^2\binom nk.$$

Първото неравенство е в сила, защото $P\to C(P)$ е инекция от множеството на изпъкналите многоъгълници с върхове в $X$ към множеството от всички подмножества на $X$. Използвайки

$$\binom ab=\frac ab\binom{a-1}{b-1}\quad\text{и}\quad\sum_{i=0}^m\binom mi=2^m,$$

пресмятаме:

$$
\begin{aligned}
Q&=\sum_{k=3}^nk^2\binom nk=n\sum_{k=3}^nk\binom{n-1}{k-1}\\
&=n\sum_{k=3}^n(k-1)\binom{n-1}{k-1}+n\sum_{k=3}^n\binom{n-1}{k-1}\\
&=n(n-1)\sum_{k=3}^n\binom{n-2}{k-2}+n\sum_{k=3}^n\binom{n-1}{k-1}\\
&=n(n-1)(2^{n-2}-1)+n(2^{n-1}-(n-1)-1)\\
&=n(n+1)2^{n-2}-2n(n-1)-n.\tag{1}
\end{aligned}
$$

От (1) следва

$$m\leq(n+1)2^{n-2}-2n+1.$$

Да разположим сега учениците във върховете на правилен $n$-ъгълник. Тогава за всяко $A\subset X$, изпъкналата обвивка на $A$ се състои от всички върхове на $A$. Значи в (1) равенството се достига и $Q=n(n+1)2^{n-2}-n-2n(n-1)$. От друга страна поради симетрията всеки ученик получава равен брой бонбони и значи $m=Q/n=(n+1)2^{n-2}-2n+1$.

Оценяване. (7 точки) 5 т. за доказване оценката отгоре, 2 т. – че тя се достига (примера). При валидна оценка отгоре, но липса на аргументация, че $P\to C(P)$ е инекция (или еквивалентно разсъждение) се отнема 1 т. Ако формулата за $m$ не е в затворен вид, (т.е. пресмятанията в (1) не са направени) се отнема 1 т.

Задача 11.4. Да се намери най-малкото естествено число $n$, за което съществуват $n$ две по две различни естествени числа $a_1,a_2,\ldots,a_n$, такива че стойността на израза

$$\frac{(a_1+a_2+\cdots+a_n)^2-2025}{a_1^2+a_2^2+\cdots+a_n^2}$$

е естествено (т.е. цяло положително) число.

Отговор. $n=9$.

Решение. Нека означим

$$S=\sum_{i=1}^na_i;\quad Q=\sum_{i=1}^na_i^2.$$

Тъй като $a_i$ и $a_i^2$ са от еднаква четност, то $S$ и $Q$ също са от еднаква четност. Щом $Q$ дели $S^2-2025$, то $S$ и $Q$ са нечетни. Но тогава $S^2-2025=(S-45)(S+45)$ се дели на 8 и понеже $Q$ е нечетно, получаваме

$$\frac{(a_1+a_2+\cdots+a_n)^2-2025}{(a_1^2+a_2^2+\cdots+a_n^2)}\geq8.$$

От друга страна от неравенството между средно аритметично и средно квадратично имаме:

$$nQ\geq S^2>S^2-2025\geq8Q,$$

следователно $n>8$, значи $n\geq9$. Ще покажем, че за $n=9$ е възможно да удовлетворим условията. Искаме да намерим решение на уравнението $S^2-2025=8Q$, тъй като знаем, че ако изразът има стойност, различна от 8, то $n>16$. Целейки симетрия, нека положим $(a_1,a_2,a_3)=(x-1,x,x+1)$, $(a_4,a_5,a_6)=(y-1,y,y+1)$ и $(a_7,a_8,a_9)=(z-1,z,z+1)$ за естествени числа $x,y$ и $z$. Тогава можем да пренапишем уравнението $S^2-2025=8Q$ като

$$(3x+3y+3z)^2-2025=8(3x^2+3y^2+3z^2+6).$$

Разделяйки двете страни на 3 и разлагайки лявата страна, получаваме:

$$3(x+y+z-15)(x+y+z+15)=8(x^2+y^2+z^2+2).$$

От съображения по модул 3 за дясната страна, точно едно от числата $x,y$ и $z$ не е кратно на 3. Нека $x=3u+1$, $y=3v$, $z=3w$ за $u,v,w\in\mathbb N$ (избрахме $x=3u+1$, но решения могат да се намерят и в случая $x=3u+2$). Пренаписвайки уравнението отново, получаваме:

$$(3u+3v+3w-14)(3u+3v+3w+16)=8(3u^2+3v^2+3w^2+2u+1).$$

Можем да забележим, че $u\equiv2\pmod3$, разглеждайки уравнението отново по модул 3. При $u=2$ уравнението е еквивалентно на

$$3(v-w)^2+\frac12(2v-7)^2+\frac12(2w-7)^2+55=0.$$

Последното няма решения, понеже лявата страна е положителна, а при $u=5$ намираме $(v;w)=(7;10)$, което води до решението

$$(a_1,a_2,\ldots,a_9)=(15,16,17,20,21,22,29,30,31).$$

Оценяване. (7 точки) 5 т. за $n\geq9$, (1 т. за доказване $n\geq c$, където $c\in[3..8]$); 2 т. за показване, че $n=9$ е възможен (пример).
